\scanpage{55}
% archive foliation: 54
\begin{center}- 12 -\end{center}
\begin{flushright}n°385\end{flushright}
\uline{caractérisé par la propriété suivante} :

\uline{Pour tout} $\mu$ \uline{dans} $\bar D_a^1(X)$ \uline{et tout système de coordonnées} $x_1,\ldots,x_n$ \uline{au voisinage de} $a$, \uline{le vecteur tangent} $u(\mu)$ \uline{admet les composantes} $\langle\mu,x_1\rangle,\ldots,\langle\mu,x_n\rangle$ \uline{par rapport au système de coordonnées} $x_1,\ldots,x_n$.

Nous allons, pour terminer ce n°, examiner l'effet d'un morphisme sur une distribution d'ordre 1.

Soit $\varphi$ un morphisme de $X$ dans une variété $X'$ de classe $C^r$. Soit $x_1,\ldots,x_n$ un système de coordonnées au voisinage d'un point $a$ de $X$, et $x'_1,\ldots,x'_m$ un système de coordonnées sur $X'$ au voisinage du point $a'=\varphi(a)$. Le morphisme $\varphi$ s'exprime dans ces systèmes de coordonnées par les formules :
\[
(28)\qquad x'_j(\varphi(p))=\varphi_j(x_1(p),\ldots,x_n(p))
\qquad 1\leq j\leq m.
\]
D'après les formules (15) et (24), on a alors :
\[
(29)\qquad
\langle\varphi_a\mu,x'_j\rangle
=\sum_{1\leq i\leq n}\langle\mu,x_i\rangle\,D_i\varphi_j(x_1(a),\ldots,x_n(a))
\qquad 1\leq j\leq m
\]
pour toute distribution $\mu\in L$ ; le premier membre de cette formule représente la j-ième composante du vecteur tangent $u_{\varphi(a)}(\varphi_a(\mu))$ dans le système de coordonnées $x'_1,\ldots,x'_m$, et, par définition de l'application $T_a(\varphi)$, le second membre de (29) représente la j-ième composante du vecteur tangent $(T_a(\varphi)\circ u_a)(\mu)$ dans le même système de coordonnées sur $X'$. Autrement dit, on a :
\[
(30)\qquad u_{\varphi(a)}\circ\varphi_a=T_a(\varphi)\circ u_a.
\]

Nous identifierons désormais les espaces $\bar D_a^1(X)$ et $T_a(X)$ au moyen de l'isomorphisme $u_a$ précédent, lorsque $r\neq\unsure{\omega|0}$.

\uline{Remarque}.-- Lorsque $r=\unsure{\omega|0}$, l'espace $T_a(X)$ n'a pas été défini précédemment [et pour cause !] ; nous poserons par définition $T_a(X)=\bar D_a^1(X)$, et $T_a(\varphi)=\varphi_a$. La formule (27) reste valable si l'on définit la matrice jacobienne par $J_{i,j}(a)=(D_i)_a y_j$, la base $(D_i)_a$ de $\bar D_a^1(X)$ étant définie par $\langle(D_i)_a,x_j\rangle=\delta_{i,j}$.
